% M2 proof module: audited WP20 multi-gap shared-Hessian action sum.

\subsection{Multi-gap shared-Hessian spectral-action sum}
\label{supp:multigap-full-proof}

Let $\rho(\bm x,\bm y)$ be one common $C^2$ physical family with baseline $\rho_0$, support projector $P$, and kernel projector $Q$. For each positive gap $\nu_k$, let the complex tangent be $A_k$ and assume a pure-boundary decomposition
\begin{equation}
P A_k P=0,
\qquad
Q A_k Q=0.
\end{equation}
Define
\begin{equation}
X_k=Q A_k P,
\qquad
Y_k=Q A_k^\dagger P,
\end{equation}
so
\begin{equation}
A_k=X_k+Y_k^\dagger.
\end{equation}
The cosine and sine support-to-kernel blocks are
\begin{equation}
K_{k,c}=\frac{X_k+Y_k}{2},
\qquad
K_{k,s}=\frac{X_k-Y_k}{2i}.
\end{equation}

For the $x_k$ coordinate, second-order positivity gives
\begin{equation}
Q\partial_{x_k}^2\rho(0)Q
\succeq2K_{k,c}\rho_0^+K_{k,c}^\dagger.
\end{equation}
For the $y_k$ coordinate,
\begin{equation}
Q\partial_{y_k}^2\rho(0)Q
\succeq2K_{k,s}\rho_0^+K_{k,s}^\dagger.
\end{equation}
Adding the two inequalities cancels the mixed $X_k/Y_k$ terms exactly:
\begin{align}
C_{\Delta,k}
&:=Q(\partial_{x_k}^2+\partial_{y_k}^2)\rho(0)Q\\
&\succeq
X_k\rho_0^+X_k^\dagger
+Y_k\rho_0^+Y_k^\dagger\\
&=:Z_{k,+}+Z_{k,-}.
\label{eq:mode-shared-curvature-proof}
\end{align}
Now sum the operator inequalities over all mode coordinates before introducing any scalar resource:
\begin{equation}
\boxed{
C_\Sigma
:=Q\sum_k(\partial_{x_k}^2+\partial_{y_k}^2)\rho(0)Q
\succeq
\sum_k(Z_{k,+}+Z_{k,-}).
}
\label{eq:common-hessian-proof}
\end{equation}
No mixed second derivatives are needed because $C_\Sigma$ is the trace of the parameter Hessian in the chosen orthonormal coordinate system.

Let $G\succeq0$ be one positive kernel cost operator and define
\begin{equation}
\Aact_{G,\Sigma}^{(2)}
=\frac14\Tr(GC_\Sigma).
\end{equation}
For each nonzero orientation let
\begin{equation}
g_{k,+}
=\lambda_{\min}
(R_{k,+}GR_{k,+}|_{R_{k,+}}),
\end{equation}
\begin{equation}
g_{k,-}
=\lambda_{\min}
(R_{k,-}GR_{k,-}|_{R_{k,-}}),
\end{equation}
where $R_{k,\pm}=\supp Z_{k,\pm}$. From Eq.~\eqref{eq:common-hessian-proof},
\begin{equation}
4\Aact_{G,\Sigma}^{(2)}
\ge
\sum_k(g_{k,+}J_{k,+}+g_{k,-}J_{k,-}),
\label{eq:total-cost-proof}
\end{equation}
with $J_{k,\pm}=\Tr Z_{k,\pm}$.

For one fixed POVM on the $N$ copies, the bilateral boundary Minkowski law applies to each Fisher block $F_{N,k}$ of that same record:
\begin{equation}
\sqrt{\frac{\Tr F_{N,k}}{N}}
\le
\sqrt{J_{k,+}}+\sqrt{J_{k,-}}.
\end{equation}
For positive bilateral prices, define the harmonic price
\begin{equation}
\gamma_k
=\left(g_{k,+}^{-1}+g_{k,-}^{-1}\right)^{-1}.
\end{equation}
Weighted Cauchy--Schwarz gives
\begin{equation}
\gamma_k
(\sqrt{J_{k,+}}+\sqrt{J_{k,-}})^2
\le
g_{k,+}J_{k,+}+g_{k,-}J_{k,-}.
\end{equation}
For a one-sided mode, set $\gamma_k$ equal to its one nonzero orientation cost. Hence
\begin{equation}
\gamma_k\frac{\Tr F_{N,k}}{N}
\le
g_{k,+}J_{k,+}+g_{k,-}J_{k,-}.
\end{equation}
Summing and using Eq.~\eqref{eq:total-cost-proof},
\begin{equation}
\boxed{
\sum_k\gamma_k\frac{\Tr F_{N,k}}{N}
\le4\Aact_{G,\Sigma}^{(2)}.
}
\label{eq:multigap-proof-final}
\end{equation}
This theorem does not require one measurement to be individually optimal for all modes. It applies to the Fisher blocks of any one common record.

\subsubsection{Clean autonomous frequency weighting}

Suppose each mode is an exact clock--signal exchange at gap $\nu_k$ and the clean upper/lower endpoint ranges are mode separated. Choose $G$ to assign the positive combined clock-plus-signal incidence price
\begin{equation}
2\hbar\nu_k
\end{equation}
to either orientation of mode $k$. Then
\begin{equation}
g_{k,+}=g_{k,-}=2\hbar\nu_k,
\end{equation}
so
\begin{equation}
\gamma_k=\hbar\nu_k.
\end{equation}
Equation~\eqref{eq:multigap-proof-final} becomes
\begin{equation}
\boxed{
\Aact_{G,\Sigma}^{(2)}
\ge
\sum_k\frac{\hbar\nu_k}{4}
\frac{\Tr F_{N,k}}{N}.
}
\label{eq:clean-multigap-proof-final}
\end{equation}

\subsubsection{One common sharp Fourier measurement}

Take the fixed-total-excitation shell
\begin{equation}
|n\rangle
\equiv|m-n\rangle_C|m+n\rangle_S,
\qquad n=-m,\ldots,m,
\end{equation}
and baseline $|0\rangle$. For $k=1,\ldots,m$ choose
\begin{equation}
A_k=c_k(|k\rangle\langle0|+|0\rangle\langle-k|),
\qquad
\nu_k=k\omega_0.
\end{equation}
The exact common family is
\begin{align}
|\psi(\bm x,\bm y)\rangle={}&
\sqrt{1-\frac12\sum_{k=1}^m c_k^2(x_k^2+y_k^2)}|0\rangle\\
&+\sum_{k=1}^m\frac{c_k}{2}(x_k+i y_k)|-k\rangle
+\sum_{k=1}^m\frac{c_k}{2}(x_k-i y_k)|k\rangle.
\end{align}
It is physical on the open ellipsoid
\begin{equation}
\frac12\sum_kc_k^2(x_k^2+y_k^2)<1
\end{equation}
and lies entirely in one total-energy eigenspace. Its kernel Hessian is
\begin{equation}
C_\Sigma
=\sum_{k=1}^m c_k^2
\left(|-k\rangle\langle-k|+|k\rangle\langle k|\right).
\end{equation}

Let $d=2m+1$ and measure the discrete Fourier basis
\begin{equation}
|v_j\rangle
=\frac1{\sqrt d}\sum_{n=-m}^{m}e^{in\phi_j}|n\rangle,
\qquad
\phi_j=\frac{2\pi j}{d}.
\end{equation}
At the baseline $p_j=1/d$. For mode $k$,
\begin{equation}
z_{j,k}
=\langle v_j|A_k|v_j\rangle
=\frac{2c_k}{d}e^{-ik\phi_j}.
\end{equation}
Therefore the cosine and sine score components are discrete harmonics. Since $1\le k,l\le m$ and $d=2m+1$,
\begin{align}
\sum_{j=0}^{d-1}\cos(k\phi_j)\cos(l\phi_j)
&=\frac d2\delta_{kl},\\
\sum_{j=0}^{d-1}\sin(k\phi_j)\sin(l\phi_j)
&=\frac d2\delta_{kl},\\
\sum_{j=0}^{d-1}\cos(k\phi_j)\sin(l\phi_j)
&=0.
\end{align}
Thus the full common-record Fisher matrix is
\begin{equation}
F_{x_kx_l}=2c_k^2\delta_{kl},
\qquad
F_{y_ky_l}=2c_k^2\delta_{kl},
\qquad
F_{x_ky_l}=0.
\end{equation}
In particular,
\begin{equation}
\Tr F_{1,k}=4c_k^2
\end{equation}
for every $k$ simultaneously.

Choose
\begin{equation}
G=2\hbar\omega_0
\sum_{n\ne0}|n|\,|n\rangle\langle n|.
\end{equation}
Then
\begin{equation}
\Aact_{G,\Sigma}^{(2)}
=\sum_{k=1}^m\hbar\nu_k c_k^2
\end{equation}
while
\begin{equation}
\sum_{k=1}^m
\frac{\hbar\nu_k}{4}\Tr F_{1,k}
=\sum_{k=1}^m\hbar\nu_k c_k^2.
\end{equation}
Hence one common measurement simultaneously saturates every mode and the complete weighted sum.

For overlapping ranges in a general model, Eq.~\eqref{eq:multigap-proof-final} remains valid for any supplied $G\succeq0$, but the existence of a simple frequency-diagonal physical interpretation becomes an operator-design question. No uniqueness claim is made outside the clean mode-separated geometry.
